% Generated draft: requires mathematical and source review.
\documentclass[11pt]{article}
\usepackage[margin=25mm]{geometry}
\usepackage{fontspec}
\setmainfont{Latin Modern Roman}
\usepackage{amsmath,amssymb,mathtools,unicode-math}
\setmathfont{Latin Modern Math}
\usepackage{enumitem,needspace,xurl,hyperref}
\setlength{\parindent}{0pt}
\setlength{\parskip}{5pt}
\setlength{\emergencystretch}{3em}
\newcommand{\N}{\mathbb{N}}
\newcommand{\Z}{\mathbb{Z}}
\newcommand{\Q}{\mathbb{Q}}
\newcommand{\R}{\mathbb{R}}
\newcommand{\C}{\mathbb{C}}
\newcommand{\F}{\mathbb{F}}
\newcommand{\A}{\mathbb{A}}
\newcommand{\PP}{\mathbb{P}}
\newcommand{\E}{\mathbb{E}}
\newcommand{\eps}{\varepsilon}
\newcommand{\dd}{\,\mathrm{d}}
\begin{document}
\section*{UnsolvedMath ID: 2200003}
\textbf{Source problem code:} catalogue entry
\subsection{Definitions and mathematical statement}
Let $N\geq1$, let $(x_i,y_i)\in\R^2$, let $\xi_i\in\R$ for $1\leq i\leq N$, and let $\alpha\geq\tfrac12$. No distinctness of the points or sign restriction on the charges is stated. Define
\[
 q_i(x)=(x-x_i)^2+y_i^2,\qquad
 D=\{x\in\R:q_i(x)>0\text{ for all }i\},\qquad
 \Psi(x)=\sum_{i=1}^{N}\frac{\xi_i}{q_i(x)^\alpha}\quad(x\in D).
\]
For $x\in D$ the real powers are well defined. If all $y_i\ne0$, then $D=\R$; otherwise the source's formula has possible singularities on the real line and no value there is supplied. For arbitrary real $\alpha$ this function is generally not rational, despite that word in the source.

\textbf{Recorded conjecture (A. Gabrielov, D. Novikov, B. Sh.).} The function $\Psi$ has at most $N$ local maxima on the whole real line. More explicitly, once the intended domain and notion of local maximum are fixed, the proposed bound is
\[
 \#\{x:x\text{ is an intended local maximum of }\Psi\}\leq N.
\]
An ordinary non-strict local maximum at $x\in D$ means $\Psi(x)\geq\Psi(z)$ for every $z\in D$ sufficiently near $x$; a strict local maximum requires $>$ for $z\ne x$. The source does not select between these notions or restrict to isolated maxima. Since arbitrary real charges include the all-zero choice, the non-strict point-count interpretation makes every point of $D$ a maximum in that permitted degenerate case. Thus a nondegeneracy or counting convention is essential; none is silently added.

The source appends the appraisal ``seems good, but no progress.'' The supplied research summary reports the conjecture still open and even the case $N=3$, $\alpha=1$, $\xi_1=\xi_2=\xi_3=1$ unresolved, and relates it to Maxwell's conjecture and Gabrielov--Novikov--Shapiro critical-point bounds. These are supplied claims requiring verification, not a resolution of the domain and degeneracy issues above.
\subsection{Short English statement}
For arbitrary real planar charges and $\alpha\geq1/2$, bound the number of local maxima of the sum of inverse powers of squared distances along the real axis by the number of charges. The source leaves real-axis poles and degenerate or non-strict maxima unresolved, and calls the expression rational although general real exponents need not be rational.
\subsection{Sources}
\begin{itemize}
\item Local UnsolvedMath registry, permanent ID 2200003, source code catalogue entry, ``Problems Around Polynomials --- Conjecture 3''; source attribution A. Gabrielov, D. Novikov, B. Sh., with ``seems good, but no progress.'' Supplied status open, classification OPEN-TRIAGE, updated 2026-08-06.
\item Shapiro, \emph{Problems Around Polynomials: The Good, The Bad and The Ugly} (2015), source-order item 3, Conjecture 3 as printed, \url{https://arxiv.org/abs/1503.05295}. The record reports source-TeX extraction and access on 2026-07-29. It says the source presents an open item and current status requires release review. Rights: NEEDS\_REVIEW; public arXiv source TeX inspected according to the record.
\item The supplied research summary mentions Maxwell's conjecture and Gabrielov--Novikov--Shapiro critical-point bounds but supplies no specific publications or theorem references for its status assessment.
\end{itemize}
\end{document}
